\chapter{按树突数目划分的神经元}
\chaptermark{树突数目}
\label{sec:dendrites}

\section{输入数目作为电路参数}

在第~\ref{sec:neurontypes}~章中，我们按照输出端释放的物质对神经元分类；在第~\ref{sec:eetypes}~章中，按照它们作为哪种器件工作分类。还有第三个特征，工程师会立刻注意到：\emph{神经元有多少个输入}。输入由树突收集——树突是供其他神经元的轴突附着的分支（第~\ref{sec:dofa}~章，第~\ref{sec:weightstore}~节）。有些神经元根本没有树突，有些有一两个，还有一些有一整棵树，收集十万个输入。对电子工程师来说，这就是输入端的扇入（fan-in），而它几乎总是与任务相称。

树突的数目和大小为什么如此重要，从一个简单的估计就能看出。膜是一个电容器，比电容 $c_m\approx1$~\textmu F/cm$^2$；神经元连同树突的面积 $A$ 越大，要把电压升高 $\Delta V$ 所需送入的电荷就越多。不计泄漏时，输入电流 $I$ 完成这件事所需的时间为
\begin{empheq}[box=\keybox]{equation}
t \;\approx\; \frac{c_m\,A\,\Delta V}{I} .
\label{eq:dendriteload}
\end{empheq}
带几个短树突的小神经元，面积约为 $2\cdot10^{-5}$~cm$^2$，电容约 $20$~pF。一个电流为几纳安的大突触不到十分之一毫秒就能把它升高 $20\mV$。拥有大树突树的神经元面积约大十五倍，电容约 $300$~pF；一个电流约 $0.1$~nA的普通突触需要 $60\ms$ 才能给它充电——远长于泄漏时间常数，也就是说永远充不上。这样的神经元需要几十个同时到达的输入。这里的数字只是数量级，但结论并不依赖于它们：\emph{输入少——快而精确；输入多——慢，靠求和}。

\begin{figure}[t]
\centering
\begin{tikzpicture}[>=Stealth,
  soma/.style={circle,draw,very thick,fill=blue!6,minimum size=0.55cm,inner sep=0pt},
  ax/.style={->,very thick,red!70!black},
  den/.style={very thick,blue!60!black}]
% 0 dendrites: sensor on a line
\begin{scope}[xshift=0cm]
\draw[den,decorate,decoration={snake,amplitude=1pt,segment length=4pt}] (-0.9,0) -- (0,0);
\node[font=\scriptsize] at (-0.9,0.3) {传感器};
\draw[ax] (0,0) -- (1.0,0);
\draw[thick] (0.1,0) -- (0.1,0.45);
\node[soma,minimum size=0.4cm] at (0.1,0.65) {};
\node[font=\scriptsize,align=center] at (0.05,-0.75) {$0$\\带传感器的线路};
\end{scope}
% 1 dendrite
\begin{scope}[xshift=3.3cm]
\draw[den] (-0.9,0) -- (-0.28,0);
\node[soma] at (0,0) {};
\draw[ax] (0.28,0) -- (0.9,0);
\node[font=\scriptsize,align=center] at (0,-0.75) {$1$\\缓冲器，跟随器};
\end{scope}
% 2 dendrites: one per ear
\begin{scope}[xshift=6.2cm]
\draw[den] (-0.28,0) -- (-0.9,0); \draw[den] (0.28,0) -- (0.9,0);
\node[soma] at (0,0) {};
\draw[ax] (0,-0.28) -- (0,-0.6);
\node[font=\scriptsize] at (-0.9,0.25) {左}; \node[font=\scriptsize] at (0.9,0.25) {右};
\node[font=\scriptsize,align=center] at (0,-1.0) {$2$\\时间上的“与”};
\end{scope}
% ~4 short dendrites: mixer
\begin{scope}[xshift=9.0cm]
\foreach \a in {45,135,225,315}{\draw[den] (\a:0.28) -- (\a:0.7); \fill[blue!60!black] (\a:0.72) circle (0.06);}
\node[soma] at (0,0) {};
\draw[ax] (0.28,0) -- (1.0,0);
\node[font=\scriptsize,align=center] at (0,-1.0) {$\sim4$\\混合器};
\end{scope}
% many: summing tree
\begin{scope}[xshift=12.0cm]
\foreach \a in {60,90,120}{\draw[den] (0,0.28) -- ++(\a:0.55) coordinate (b); \foreach \d in {-20,20}{\draw[den] (b) -- ++(\a+\d:0.35);}}
\foreach \a in {-120,-60}{\draw[den] (0,-0.28) -- ++(\a:0.45);}
\node[soma] at (0,0) {};
\draw[ax] (0.28,0) -- (1.0,0);
\node[font=\scriptsize,align=center] at (0,-1.0) {$10^4$--$10^5$\\求和器};
\end{scope}
\end{tikzpicture}
\caption{按输入数目划分的五类。蓝色为树突（输入），红色为轴突（输出）。输入越少，神经元传递单个信号越快越精确；输入越多，它越擅长求和与分类。这是示意图，不是细胞形态图。}
\label{fig:dendriteclasses}
\end{figure}

\section{零个树突：线路末端的传感器}

触觉、痛觉和肌肉牵张的感觉神经元（第~\ref{sec:sensors}、\ref{sec:pain}、\ref{sec:muscles}~章），细胞体上没有一个树突。轴突以一根细枝离开胞体，随即一分为二：一支通往皮肤或肌肉，其末端本身就是传感器；另一支进入脊髓。脉冲从传感器直接跑进脊髓，绕过细胞体。对工程师来说，这是一条传感器位于远端的通信线路，而细胞体挂在一旁，就像电源：它为线路供能，但不参与传输。好处是速度和可靠性：路径上没有任何一处需要对信号求和、再重新决定是否传递。

光和声的传感器是更极端的情况：它们不是收集输入的神经元，而是换能器本身，其输入不是突触，而是受体蛋白（图~\ref{fig:transducer}）。

\section{一个树突：缓冲器}

有一个树突和一个轴突的神经元接收一条输入线路并把它传下去。嗅觉神经元就是这样：唯一的树突伸到鼻腔表面，只携带一种类型的受体\cite{Mombaerts2004}；视网膜中位于光传感器和眼睛输出神经元之间的中继细胞也是这样；连接耳朵传感器和大脑的神经元同样如此。工程师会称它们为缓冲器或跟随器：它们不计算，而是递送一个信号，有时会改变其形式——例如像第~\ref{sec:sensors}~章那样，把传感器的平滑电压变成脉冲。

\section{两个树突：时间上的“与”}

哺乳动物中确定声音方向的符合检测器（图~\ref{fig:itd}）常常恰好有两个树突，朝向相反：一个接收来自左耳的输入，另一个接收来自右耳的输入\cite{Grothe2010}。为什么要把输入分开？回想公式~\eqref{eq:shunt}：当膜的某一处打开了许多兴奋性通道时，那里的电压会接近平衡电位，之后每个输入增加的量越来越少。因此，来自同一只耳朵、汇集在同一个树突上的许多输入会使它饱和，单靠它们无法把神经元推到阈值。而两个树突各来一个输入，则几乎线性相加。两个树突把神经元变成更严格的“与”：只有信号从两侧同时到来时才发放。模型表明，这种分开确实改善了符合检测\cite{AgmonSnir1998}。

\section{几个短树突：混合器与中继器}

\emph{混合器。}在第~\ref{sec:motorlearning}~章的运动学习回路中，约七百亿个小\nt{GLUT}神经元把输入扩展成一个很长的向量。它们每个只有约四个短树突，每个树突末端是一个“爪”，抱住一个输入的末梢。因此，每个混合器只看到众多输入中的约四个，作为其四元组上的一个小“与”元件工作。从 $M$ 条输入线路中可以选出 $\binom{M}{4}$ 个不同的四元组：$M=100$ 时接近四百万。为什么需要这么多？一个有 $n$ 个输入的阈值元件，能正确分为两类的随机模式不超过
\begin{empheq}[box=\keybox]{equation}
P_{\max} \;\approx\; 2n
\label{eq:cover}
\end{empheq}
个\cite{Cover1965}。同一回路的输出神经元从混合器接收约 $10^5$ 个输入，按~\eqref{eq:cover}，它的上界约为 $2\cdot10^5$ 个不同的对应关系。这是理想元件的极限：如果所有权重都为正，如兴奋性突触那样，并且需要留出抗噪余量，那么对同一个神经元的估计约为 $4\cdot10^4$ 个对应关系\cite{Brunel2004}。输入少的混合器，正是为了让大求和器拥有许多不同的、几乎独立的输入\cite{Marr1969,Albus1971}。

\emph{中继器。}在从耳朵通往方向检测器的路径上，有一些只带少数短树突的神经元，它们只接收几个巨大的突触，有些实际上只接收一个。按~\eqref{eq:dendriteload}，这正是所需要的：电容小、电流大，神经元在每个输入脉冲之后不到一毫秒就发放，几乎不损失时间精度。其中一种中继器接收\nt{GLUT}，输出\nt{GLYC}：这是一个精度达几十微秒的反相器，为方向检测器提供快速抑制\cite{BorstSoria2012}。

\section{成千上万个输入：求和器与分类器}

尺度的另一端是皮层中有上万个输入的\nt{GLUT}神经元，以及运动学习回路中有十万个输入的输出\nt{GABA}神经元。按~\eqref{eq:dendriteload}，这样的神经元很慢，不能对单个输入作出响应：它做求和。这就是第~\ref{sec:glutcomputer}~章的积分神经元，也是构建分类器所用的求和器。大树突树还带来新东西：它的分支像各自带阈值的独立子电路一样工作，因此一个神经元的行为就像一个小型多层网络\cite{Beniaguev2021,Gidon2020}。

\section{同时也是输出的树突}

有些神经元根本没有轴突，树突既是输入又是输出：它们接收信号，自己也释放神经递质。研究最多的例子是视网膜中参与确定运动方向的\nt{GABA}神经元（第~\ref{sec:gaba}~章，图~\ref{fig:direction}）。这种神经元的每个分支各自对从细胞体向其末梢方向的运动作出响应，并从这个末梢输出抑制\cite{Euler2002}。一个神经元就是几个独立的方向检测器，每个分支一个。对工程师来说，这不是一个元件，而是一个封装内有多个通道的芯片。

\section{小结}

\begin{table}[t]
\centering
\footnotesize
\setlength{\tabcolsep}{3pt}
\renewcommand{\arraystretch}{1.15}
\begin{tabular}{>{\raggedright}p{1.9cm}>{\raggedright}p{3.4cm}>{\raggedright}p{2.6cm}>{\raggedright}p{1.8cm}>{\raggedright\arraybackslash}p{3.8cm}}
\toprule
树突数 & 例子 & 器件 & 输入数 & 已知什么 \\
\midrule
$0$ & 触觉、痛觉、牵张传感器 & 末端带传感器的线路 & --- & 已确定 \\
$1$ & 嗅觉神经元、视网膜中继细胞、耳的神经元 & 缓冲器，跟随器 & $1$ 到少数 & 已确定 \\
$2$ & 声音方向检测器 & 时间上的“与” & 两束 & 结构已确定；分开输入的好处已在模型中证明 \\
$\sim4$ & 运动学习回路的混合器 & 小“与”元件 & $\sim4$ & 已确定；扩展输入的作用是有充分依据的假说 \\
少数，短 & 耳通路上的中继器 & 中继器，反相器 & $1$ 到几个巨大突触 & 已确定 \\
成千上万 & 皮层\nt{GLUT}神经元，运动学习的输出神经元 & 求和器，分类器 & $10^4$--$10^5$ & 已确定；分支作为子电路已在个别例子中测量 \\
树突即输出 & 视网膜方向\nt{GABA}神经元 & 多通道芯片 & 许多 & 已测量 \\
\bottomrule
\end{tabular}
\caption{按树突数目划分的神经元。}
\label{tab:dendrites}
\end{table}

\begin{keyblock}
\begin{proposition}[输入数目跟随任务]
\label{prop:dendrites}
凡是需要单个信号的速度和精度之处——在传感器、中继器和符合检测器上——神经元的树突少、输入少、突触大：小电容~\eqref{eq:dendriteload}~充电快。凡是需要求和与分类之处，神经元都有带成千上万个输入的大树突树。各类的结构已有测量；计算上的解释有充分依据，但对许多类别来说仍是假说。
\end{proposition}
\end{keyblock}

表~\ref{tab:dendrites}~补充了前两种分类：按神经递质（表~\ref{tab:types}）和按器件（表~\ref{tab:eetypes}）。三种分类都是从不同侧面观察同一个元件：它输出什么，它计算什么，以及它听多少。

\section{习题}

\begin{exercise}[\lvlA]
\label{ex:19:1}
\emph{大神经元需要多少输入。}一个 $C=300$~pF、$\tau_m=20\ms$ 的神经元接收突触，每个突触是 $0.1$~nA的电流源，阈值比静息高 $20\mV$。(a)~要能到达阈值，至少需要多少个突触同时工作？在 $10\ms$ 内呢？$5\ms$ 内呢？(b)~$C=20$~pF的神经元在一个 $4$~nA突触作用下多久到达阈值？
\end{exercise}

\begin{exercise}[\lvlA]
\label{ex:19:2}
\emph{为什么是四个。}混合器是对 $M=100$ 条线路中 $k$ 条的“与”，每个模式中一半线路活跃，混合器共 $7\cdot10^{10}$ 个。(a)~$k=2$、$4$、$40$ 时各有多少个混合器对一个模式响应？(b)~与 $M=100$ 条直接线路相比，混合器使有 $10^5$ 个输入的输出神经元的对应关系数~\eqref{eq:cover}~增加了多少倍？
\end{exercise}

\begin{exercise}[\lvlB]
\label{ex:19:3}
\emph{$2n$ 从何而来。}过原点的超平面把 $n$ 维空间中处于一般位置的 $P$ 个点分成两类，共有 $C(P,n)=2\sum_{k=0}^{n-1}\binom{P-1}{k}$ 种方式。(a)~证明当 $P=2n$ 时，$2^P$ 种划分中恰好一半可分。(b)~$n=100$、$P=180$ 和 $220$ 时，可分的比例是多少？
\end{exercise}

\begin{exercise}[\lvlB]
\label{ex:19:4}
\emph{两个树突作为严格的“与”。}树突是泄漏电导为 $g_L$ 的一段，每个输入在其上打开电导 $g_0=0.5\,g_L$，反转电位为 $E_E$，胞体看到 $\kappa(V_1+V_2)$，阈值 $\theta=1.1\,\kappa E_E$。为什么同一个树突上无论多少输入都无法触发神经元？每个树突上需要多少输入？
\end{exercise}

\begin{exercise}[\lvlB]
\label{ex:19:5}
\emph{设计中继器。}脉冲抖动 $\sigma_t=\sigma_V/\dot V$，其中 $\dot V$ 是阈值附近的电压斜率。要求在噪声 $\sigma_V=1\mV$ 下 $\sigma_t\le20$~\textmu s；忽略泄漏。$C=300$~pF的神经元为此需要多少个 $0.1$~nA的同步突触？$C=20$~pF、只有一个 $4$~nA突触的神经元抖动是多少？
\end{exercise}

\begin{exercise}[\lvlB]
\label{ex:19:6}
\emph{建模：长树突的充电。}模拟长度 $L=\ell/\lambda$、两端封闭的无源电缆 $\tau_m\partial_tV=\lambda^2\partial_x^2V-V+i/g$（$40$ 段，欧拉法）。在远端注入短电流脉冲：$L=0.2$、$1$、$2$ 时，近端何时升到、升到多高（以同样电荷均摊到整个树突时电压的比例计）？
\end{exercise}

\begin{exercise}[\lvlC]
\label{ex:19:7}
\emph{研究：最优输入数目。}电容 $C=C_0+nc_s$，同时有 $m$ 个输入工作，每个电流 $I_s$，神经元记忆 $P\approx2n$ 个对应关系~\eqref{eq:cover}。在 $m$ 恒定和比例 $m=\varphi n$ 恒定时，响应时间如何依赖于 $P$？提出一个代价准则，并分别求中继器和分类器的最优 $n$。
\end{exercise}
